% Build: latexmk -pdf main.tex
% Draft. Figures are deliberately absent: every figure is a \figph placeholder
% stating what the panel must show and which pipeline output produces it.

\documentclass[11pt,a4paper]{article}

\usepackage[T1]{fontenc}
\usepackage[utf8]{inputenc}
\usepackage[english]{babel}
\usepackage{amsmath,amssymb,amsthm,mathtools}
\usepackage{microtype}
\usepackage{geometry}
\usepackage{booktabs}
\usepackage[dvipsnames]{xcolor}
\usepackage[round]{natbib}
%\usepackage{mathptmx}
\usepackage{hyperref}

\geometry{margin=2.5cm}
\hypersetup{colorlinks=true, allcolors=RoyalBlue}

% Custom commands and theorem environments.

\newcommand{\N}{\mathbb{N}}
\newcommand{\Z}{\mathbb{Z}}
\newcommand{\R}{\mathbb{R}}

\theoremstyle{plain}
\newtheorem{theorem}{Theorem}[section]
\newtheorem{lemma}[theorem]{Lemma}
\newtheorem{proposition}[theorem]{Proposition}
\newtheorem{corollary}[theorem]{Corollary}

\theoremstyle{definition}
\newtheorem{definition}[theorem]{Definition}
\newtheorem{example}[theorem]{Example}

\theoremstyle{remark}
\newtheorem{remark}[theorem]{Remark}


\title{The shape of attention:\\
detecting and classifying anomalous lexical activity\\
in eighty years of the daily press}

\author{%
  Elias Echikr \and
  Benoît de Courson \and
  Simon Coste
}

\date{\today}

\begin{document}
\maketitle

\begin{abstract}blabla
\end{abstract}

\tableofcontents

\section{Introduction (pur AI slop à récrire)}

Public attention behaves like an excitable medium. Most of the time a given
topic is quiescent; occasionally, it is triggered --- by an event, by a
rumour, by its own internal dynamics --- and then dominates the available
bandwidth for hours, days, or years before relaxing. This qualitative picture
has been made quantitative in a series of studies of digital traces:
best-seller rankings \citep{sornette2004endogenous}, video views
\citep{crane2008robust}, news-story lifetimes \citep{wu2007novelty}, hashtag
volumes \citep{lehmann2012dynamical}. Two ideas recur. The first is that
bursts are not exceptions to be discarded but the observable of interest, and
that their statistics are heavy-tailed
\citep{barabasi2005origin,goh2008burstiness}. The second is that the
\emph{shape} of a burst is informative: an externally imposed shock produces a
different response function than a shock generated by the system's own
self-excitation, so that shapes can in principle be used to classify
mechanisms \citep{sornette2004endogenous,crane2008robust}, in the same spirit
as reflexivity measurements in finance
\citep{filimonov2012quantifying,hardiman2013critical}.

Digital traces, however, are awkward measuring devices. They rarely extend
beyond a decade, their sampling and retention policies change under the
researcher's feet, and --- most damaging --- they usually lack a denominator:
a rise in the absolute number of posts mentioning a word may reflect nothing
but a rise in the number of posts. The daily press has the opposite profile.
A newspaper published every day for eighty years is a metronome. The number of
words it prints each day is measurable, which gives an exact exposure. Its
archives are stable. And because several titles cover the same news with
different editorial priorities, the same event is recorded several times by
partly independent instruments --- a property no single social platform
offers, and the one on which any causal reading of attention dynamics
ultimately rests.

What the press lacks, so far, is a statistical theory of its own noise. Daily
word counts are integers, small for most words, strongly overdispersed, and
often zero; the corpus size fluctuates by more than two orders of magnitude
over eight decades. Existing work on lexical time series either smooths these
difficulties away by working with frequencies and eyeballing curves
\citep{michel2011quantitative,azoulay2023gallicagram}, or detects bursts with
model-free heuristics \citep{kleinberg2002bursty} whose output cannot be
compared across words of different frequency. Yet the statistical
characterisation of word counts has been available for thirty years: content
words are bursty, and mixtures of Poisson laws --- the negative binomial being
the canonical case --- describe them well \citep{church1995poisson,katz1996distribution}.
What has been missing is the step from that description to a calibrated
detector: a null model, fitted word by word with the correct exposure, under
which each day receives a $p$-value, and every anomalous day in a corpus of
$10^{4}$ words over $10^{4}$ days becomes comparable to every other on a
single scale.

This paper takes that step and then uses it. Our contributions are the
following.

\begin{enumerate}
  \item \textbf{A day-resolved lexical corpus of the French daily press}
    (Section~\ref{sec:data}), built by systematic mapping and scraping of
    online archives, with per-day $n$-gram counts and per-day total word
    counts for four titles, the longest of which covers \emph{Le Monde} from
    December 1944 to 2025.
  \item \textbf{A calibrated null model for daily word counts}
    (Section~\ref{sec:detection}): a Bernoulli times shifted negative
    binomial mixture with exposure proportional to the daily corpus size,
    estimated word by word, robustified by a two-pass fit, and validated by
    per-word goodness-of-fit diagnostics. We also state precisely what the
    model does \emph{not} do --- it is a model of stationary noise, not of
    regime change --- and quantify the cost of that limitation.
  \item \textbf{A catalogue of attention events and its statistics}
    (Section~\ref{sec:detection}): $164\,254$ anomalous (word, day) pairs on
    $9\,817$ words, reduced to $123\,465$ distinct events by a
    non-transitive greedy suppression step borrowed from seismological
    declustering and object detection.
  \item \textbf{The geometry of single-jump shapes}
    (Section~\ref{sec:single}): windows of $\pm 15$ publication days around
    each event, normalised per window, analysed by principal component
    analysis. The result is negative and, we believe, robust: the spectrum is
    flat (six components carry $33.3\%$ of the variance against $20\%$ for an
    isotropic cloud of the same rank), kernel principal component analysis
    does not concentrate it further, and the leading shapes are reproduced at
    three temporal resolutions. Attention jumps in the press form a continuum
    of amplitudes and durations, not a small taxonomy of response functions.
\end{enumerate}

Section~\ref{sec:journals} states the design we consider decisive and which
this draft does not yet deliver: comparing the same event across titles. A
jump seen simultaneously and proportionally by every newspaper is a property
of the world; a jump seen by one title alone, or earlier by one title than by
the others, is a property of that newspaper. This is the identification that
separates exogenous from endogenous excitation, and it is unavailable to any
single-corpus study. Section~\ref{sec:conclusion} concludes.

\section{Data collection}
\label{sec:data}

partie à écrire par Elias !

\section{Detecting anomalous activity}
\label{sec:detection}


Fix a word $w$. For each publication day $t \in \{1, \dots, T\}$
we observe the count $X_t \equiv X^w_t \in \N$ of occurrences of $w$, and the \emph{exposure} of that day, which is the
total number of tokens $N_t$ published that day. The word frequency is 
\[
  f_t = 10^{5}\, \frac{X_t}{N_t}.
\]
Our goal is to detect days where $f_t$ is unusually high; but we have to account for the fact that $N_t$ varies from day to day. A day with low $N_t$ is more likely to have a high $f_t$ simply because it has fewer tokens to distribute, and thus $f_t$ would have a higher variance. To detect these unusually high frequencies while taking into account $N_t$, we first fit a statistical model to the data, and then we use the model to compute $p$-values for each day. 

\subsection{Statistical model and fitting}
\label{sec:model}

We model the daily count as a mixture: a Bernoulli variable
decides whether the word is used at all that day (there are many days where the word is not used at all). Then, conditionally on being
used, the count is a negative binomial variable (minus one) whose mean is
proportional to $N_t$. Formally, with parameters
$\theta = (p_0, \mu_b, r_b) \in (0,1) \times (0,\infty)^{2}$,
\begin{equation}
  \PP_\theta(X_t = k) =
  \begin{cases}
    p_0, & k = 0,\\[0.4em]
    (1 - p_0)\, f_{\NB}\!\left(k - 1;\ \mu_b N_t,\ r_b\right), & k \geq 1,
  \end{cases}
  \label{eq:model}
\end{equation}
with $f_{\NB}$ being the density of the negative binomial, 
\begin{equation}\label{eq:nbpdf}
  f_{\NB}(k; \mu, r) = \frac{\Gamma(k + r)}{\Gamma(k + 1) \Gamma(r)} \left(\frac{\mu}{r + \mu}\right)^r \left(\frac{r}{r + \mu}\right)^k,
\end{equation}
In the sequel, we simply note $X_t \sim \BNB(p_0, \mu_b N_t, r_b)$. Conditionnally on being used, the mean of $X_t$ is proportional to $N_t$: this is standard in the litterature for regression of count data \citep{cameron2013regression,hilbe2011negative}.

Note that \eqref{eq:model} is a hurdle model \citep{cragg1971some,mullahy1986specification}: the zeros are carried \emph{only} by the Bernoulli component, and the counting component is supported on $\{1, 2, \dots\}$. In fact the two
regimes are meant to be interpreted separately: ``is the word active'' and ``how intense is the activity''  \citep{lambert1992zero}.
The moments of the model are
\begin{align}
  \E_\theta[X_t] &= (1 - p_0)\,(1 + \mu_b N_t), \\
  \Var_\theta(X_t) &= (1 - p_0)\left(\mu_b N_t + \frac{(\mu_b N_t)^{2}}{r_b}\right)
   + p_0 (1 - p_0)\,(1 + \mu_b N_t)^{2},
  \label{eq:varmix}
\end{align}
the second term of \eqref{eq:varmix} being the extra variance due to the Bernoulli component. 

\begin{remark}talk abt overdispersion\end{remark}

 


We will fit models like \eqref{eq:model} to each word in the corpus, using maximum likelihood estimation. The likelihood of \eqref{eq:model} is easy to evaluate: writing $n_0 = \#\{t : X_t = 0\}$ for the number of inactive days and
$\mathcal{A} = \{t : X_t \geq 1\}$ for the set of active days, the log-likelihood is
\begin{equation}
  \ell(\theta)
  = \underbrace{n_0 \log p_0 + (T - n_0) \log (1 - p_0)}_{\text{Bernoulli}}
  + \underbrace{\sum_{t \in \mathcal{A}}
      \log f_{\NB}\!\left(X_t - 1;\ \mu_b N_t,\ r_b\right)}_{\text{counting}},
  \label{eq:loglik}
\end{equation}
in which $p_0$ appears only in the first bracket and $(\mu_b, r_b)$ only in
the second. Hence the maximum likelihood estimator gives $\hat p_0 = n_0/T$
on the first hand, and then $(\hat\mu_b, \hat r_b)$ maximises a standard negative binomial likelihood on the active days $\mathcal{A}$, which is done using standard methods implemented in the Scipy library. 

\subsection{\texorpdfstring{$p$}{p}-values and spike detection}
\label{sec:pvalues}

Given the fitted parameters$\hat\theta$, each active day receives an upper-tail $p$-value under the
fitted law,
\begin{equation}
  p_t = \PP_{\hat\theta}(X \geq X_t \mid N_t)
      = (1 - \hat p_0)\,
        \PP\!\left(\NB(\hat\mu_b N_t, \hat r_b) \geq X_t - 1\right),
  \qquad X_t \geq 1,
  \label{eq:pvalue}
\end{equation}
and we call
\[
  \surp_t = -\log_{10} p_t
\]
the \emph{surprise} of the day. Given a threshold $H$, a day is declared anomalous\footnote{We only test for excess activity and we did not consider the opposite, an anomalous absence of activity. } when
$\surp_t \geq H$, that is $p_t < 10^{-H}$. In the sequel, we use $H=6$ to be very conservative. If the statistical fit is correct, that would correspond to a one-in-a-million event. 



\begin{remark}Two-stage fitting and robustification\end{remark}

\begin{remark}
Given fitted parameters $\hat{\theta}$, the quantiles of $\mathrm{BNB}(\hat p_0, \hat\mu_b N, \hat r_b)$ are decreasing with the exposure $N$: with a low exposure $N$, the count $X$ has a higher variance and thus a higher probability to be large. Consequently, two different days $t\neq t'$ with the same word count $X_t = X_{t'}$ can have a very different surprise, if the exposure of the two days are different, $N_t \neq N_{t'}$. Indeed, if $N_t < N_{t'}$, then $s_t < s_{t'}$.
\end{remark}

\subsection{Statistics of activity spikes}
\label{sec:catalogue}

number of spikes / surprise statistics / co-jumps and heavy tails

\section{Classification of single jumps (PUR AI SLOP)}
\label{sec:single}

The catalogue turns each attention event into a dated point. This section asks
a shape question: around an event, what does the trajectory of the word look
like, and do those trajectories fall into a small number of classes? The
motivation is the response-function programme of
\citet{sornette2004endogenous} and \citet{crane2008robust}, where the profile
of a burst is read as a signature of its mechanism, and the analogous
classification of collective attention on social media
\citep{lehmann2012dynamical}.

\subsection{Windows, and the normalisation trap}

Around each retained event we extract the frequency series on
$\pm d = 15$ publication days, a vector of $D = 2d + 1 = 31$ numbers. Two
choices in this apparently innocuous step do most of the work.

The window is measured in publication days, not calendar days, so that
holidays and strikes do not stretch it. And each window is standardised
\emph{along itself} --- its own mean subtracted, its own standard deviation
divided out --- so that only its shape remains. The alternative
standardisations are instructive failures. Raw frequencies make the analysis
discover that some words are more frequent than others. Standardising column
by column, which is what the normalisation option of standard implementations
does, does not fix this: it puts $63.5\%$ of the variance on a first component
whose projection correlates at $0.99$ with the raw mean level of the window,
that is on the direction ``is this word frequent''. We keep that variant as an
explicit negative control. A per-window rescaling to $[0,1]$ behaves better but
its first component still carries the overall level (correlation $0.49$). Only
the per-window $z$-score removes level and amplitude and leaves shape.

Per-window standardisation has one structural consequence that must be kept in
mind when reading spectra: a $z$-scored window sums to zero, so the cloud lives
in $D - 1 = 30$ dimensions and a structureless cloud would put
$1/30 \approx 3.33\%$ of its variance on each component. This is the baseline
against which every percentage below should be read.

\subsection{Thin days inside windows}

The pathological days of Figure~\ref{fig:corpussize} contaminate windows even
when they are not themselves events: a day with $N_t = 177$ tokens sends
$f_t = 10^{5} X_t / N_t$ through the roof, and $17\%$ of windows contain at
least one day below $5\,000$ tokens. The symptom was found by inspection of
the fourth component, whose extreme projections were $17.8$ times more frequent
in windows containing such a day than in clean ones. We therefore replace, in
each window, the frequencies of days below $5\,000$ tokens by linear
interpolation from their healthy neighbours, and discard the $1\,505$ windows
whose \emph{central} day is itself thin: their detection remains valid, since
the null model accounts for $N_t$, but the shape of their window does not.
The cleaning leaves $121\,805$ windows.

This is a case where a control changed the data and not the conclusions, which
is worth reporting as such. After cleaning, the enrichment of the fourth
component falls from $17.8$ to $0.5$ and its archetypal windows become
interpretable events; the variance spectrum moves by less than half a point per
component; the leading six directions are preserved with $|\cos| \geq 0.94$;
and the whole thing is insensitive to the $5\,000$-token threshold, the leading
three-dimensional subspace rotating by less than $4^{\circ}$ when the threshold
is moved between $2\,000$ and $10\,000$.

\subsection{A flat spectrum}

Principal component analysis of the $121\,805 \times 31$ matrix of cleaned,
$z$-scored windows gives the spectrum
\[
  9.2,\ 6.2,\ 5.3,\ 4.6,\ 4.2,\ 3.8\ \%,
\]
that is $33.3\%$ of the variance in six components, against $20\%$ for an
isotropic cloud of rank $30$; fifteen components are needed to reach $62\%$.
The excess over isotropy is real and highly significant given the sample size,
but it is small. There is no pair or triplet of dominant profiles: the
diversity of window shapes is genuine, and it is the first result of this
section.

\figph{spectrum}{Variance spectrum of the window cloud. The observed spectrum,
the isotropic baseline at $3.33\%$ per component, and two null spectra
obtained by destroying structure in the data (independent shuffling of each
column, and circular shift of each window).}{Reads the spectra of the
reference configuration and of its two nulls. The nulls are what make the
excess interpretable, and what makes spectra of different window widths
comparable at all.}

\figph{normalisation}{Spectra under the three normalisations: per-window
$z$-score, per-window rescaling to $[0,1]$, and the column-by-column negative
control.}{Must show the control collapsing more than $60\%$ of the variance
onto an artefactual first component.}

\subsection{The leading shapes}

The first four components are readable temporal profiles, and their archetypal
windows are recognisable dates. The first is an isolated one-day spike, dipping
slightly on the days before and after. The second is a durable elevation: the
word rises at the event and does not come back down --- an installation rather
than a spike. The third is a before/after switch centred on the event day. The
fourth is a hollow before the event followed by a rebound. Components five and
six are a pair of slow oscillations; we tested and rejected the natural
hypothesis that they encode the weekly rhythm of the paper, since their phase
is unrelated to the weekday of the event, and we read them as generic
oscillatory modes of a smooth $31$-point window rather than as a signal.

\figph{components}{The first six components as temporal profiles, aligned on
the event day.}{Reads the component matrix of the reference configuration.}

\figph{archetypes}{For each of the first four components, the three real
windows most aligned with it, labelled by word and date.}{Reads the
projections and the window metadata. The purpose is to let the reader verify
the interpretation of the profiles against dates they can recognise.}

\figph{plane12}{The window cloud in the plane of the first two components
(density, log scale), with the most surprising window of a few landmark event
words marked.}{Must show that the interpretation of the axes --- isolated
spike against durable installation --- separates episodes the reader knows.}

\subsection{A continuum, not a taxonomy}

A flat spectrum is consistent with two very different situations: a genuinely
featureless continuum of shapes, or a low-dimensional but \emph{curved} family
that a linear method cannot compress. We tested the second explicitly with
kernel principal component analysis \citep{scholkopf1998nonlinear} using a
Gaussian kernel, with the bandwidth calibrated from the median pairwise
distance of the cloud and then scanned over two decades around it. It never
concentrates the variance better than the linear analysis: on a sample of
$8\,000$ windows of the reference configuration, the six leading components
carry $33.0\%$ of the variance linearly against $24.9\%$ with the kernel at
$\gamma = 1/D$, and the same comparison on \emph{Les Échos} gives $34.3\%$
against $27.2\%$. (This
comparison is only meaningful once the variance is normalised by the trace of
the centred Gram matrix; with a truncated denominator the kernel appears
spuriously favourable.) The absence of a curved structure supports the direct
reading: the windows form a continuum of amplitudes and widths around a single
peaked profile, so the components describe deformations of one shape rather
than a set of distinct shapes.

The second test is scale. The entire pipeline was rerun on grids of $3$ and $7$
consecutive publication days, with counts and exposures summed within blocks
and all downstream settings unchanged. The leading profiles survive: cosines
with their daily homologues are $0.98$, $0.92$, $0.90$ on the first three
components for the three-day grid and $0.97$, $0.96$, $0.98$ for the weekly
one. The spectrum concentrates only slightly as the grid coarsens, from
$33.3\%$ to $37.9\%$ and $41.2\%$ at six components. Isolated spike, durable
installation and before/after switch therefore exist at the scale of the day
and at the scale of the week alike; they are not artefacts of daily sampling.
Aggregation is not free, however: it acts as a low-pass filter on detection
itself, and $2\,800$ words that have events on the daily grid have none on the
weekly one, so the three grids do not observe the same population of events.

\figph{kernel}{Linear against kernel principal component analysis. Cumulative
variance of the leading components for the linear analysis and for the
Gaussian kernel at three bandwidths, on two titles.}{Must make clear that the
comparison uses the trace of the centred Gram matrix as denominator.}

\figph{scales}{The leading four components on the daily, three-day and weekly
grids, and the matrix of cosines between homologous components.}{Reads the
three grid configurations. The point is the stability of the profiles, and the
loss of the $2\,800$ words that no longer have detections after aggregation.}

\subsection{Comparing configurations honestly}

The results above depend on choices --- window width, temporal grid, detection
threshold, vocabulary, transformation --- and comparing configurations by their
variance spectra is a trap that we fell into before instrumenting against it.
Three effects must be neutralised. Wider windows have more components and
mechanically flatter spectra, so spectra of different $D$ are not comparable.
Larger samples give more stable but not more concentrated spectra, so
configurations must be compared at matched sample size --- we resample to
$5\,000$ windows with several seeds, and measure a between-seed dispersion of
$0.0045$ on the summary statistics. And a spectrum concentrates for reasons
that have nothing to do with the event: any smooth curve has a compressible
spectrum. We therefore report each spectrum against two nulls --- independent
shuffling of each column, which destroys temporal structure entirely, and
circular shifting of each window, which preserves it but destroys its
alignment on the event day. The ratio to the second null is the only statistic
that measures what we care about: the part of the structure that is anchored on
the jump.

Under these controls, a campaign over roughly four thousand configurations
yields two robust statements and one negative one. A higher detection threshold
buys a much better anchored structure ($1.65$ against $1.31$ at threshold
$\surp \geq 6$ and $4$ respectively). A logarithmic transformation of the
frequencies surfaces slow texture, and destroys anchoring almost completely, so
we do not use it. And titles differ in a way that anticipates
Section~\ref{sec:journals}: anchoring is highest for \emph{Mediapart} ($1.62$)
and lowest for \emph{Le Figaro} ($1.28$), while excess texture runs the other
way, which is what one expects of a small pure-online title covering sharp
events against a thick generalist daily.

For completeness, the variance-optimal configuration is not the one we used
above: on \emph{Le Monde} with weekly blocks, $\pm 10$ blocks and threshold
$\surp \geq 6$, six components carry $60.4\%$ of the variance and the first
alone $18.5\%$, four components suffice to reach half the variance, and six
reconstruct a typical window to $96\%$. It achieves this on $8\,764$ windows,
fourteen times fewer than the reference configuration, and with an anchoring
ratio no better. Reporting it as ``the'' result would be exactly the kind of
selection our nulls were built to expose; we report it as the end of the
trade-off between compression and coverage.

\subsection{What the continuum means}

If jump shapes formed a small taxonomy, the natural next step would be the one
the response-function literature takes: read each class as a mechanism, and in
particular separate externally triggered episodes, which begin abruptly, from
endogenously grown ones, which are preceded by a slow rise
\citep{sornette2004endogenous,crane2008robust}. Our result blocks that route in
its simple form. A continuum of amplitudes and widths around one peaked profile
means that shape alone, measured on a single title, does not sort events into
mechanisms.

Two ways forward remain, and we distinguish them clearly because only the
second is available to us. The first stays within one corpus and tests the
predictions of self-excitation quantitatively rather than qualitatively:
power-law relaxation after the peak, of the Omori type
\citep{utsu1995centenary}, the presence or absence of precursory growth, and a
direct fit of a self-exciting point process to the arrival times of anomalous
days \citep{hawkes1971spectra,filimonov2012quantifying}. \TODO{decide whether
to include the relaxation-law analysis in this paper; it requires fitting decay
exponents on the event windows, which is a small addition to the existing
pipeline.} The second is comparative, and is the subject of the next section.

\section{The structure of jumps across journals}
\label{sec:journals}

totalement à écrire 

PCA sur les windows média par média

\section{Conclusion}
\label{sec:conclusion}


\section*{Data and code availability}

\TODO{state what is released: derived count series, the event catalogue, the
window matrices, the analysis code, and the public query interface. The article
corpus itself cannot be redistributed.}

\section*{Acknowledgements}

\TODO{acknowledgements, funding, and compute.}

\bibliographystyle{plainnat}
\bibliography{refs}

\end{document}
